\documentclass[11pt]{article}
\usepackage[margin=2.6cm]{geometry}
\usepackage{booktabs}
\usepackage{graphicx}
\usepackage{siunitx}
\usepackage{amsmath}
\usepackage{hyperref}
\usepackage{seqsplit}
\graphicspath{{figures/}}
\newcommand{\runid}[1]{\texttt{\small\seqsplit{#1}}}
\newcommand{\dsid}[1]{\texttt{\small #1}}
\newcommand{\best}[1]{\textbf{#1}}
\newcommand{\medape}{medAPE}
\newcommand{\rsq}{$R^2$}

\title{Spotify Engagement from Metadata: Experiment Synthesis}
\author{}
\date{\today}

\begin{document}
\maketitle

\begin{abstract}
This document is the living synthesis of the taste-fit prediction campaigns;
the markdown reports in \texttt{experiments/} are the primary record. The
active target is a binary \emph{rotation} signal --- did a track earn a repeat
listen, $\text{play\_count}\ge 2$ --- predicted from a track's \emph{intrinsic
metadata} alone (genre, artist, artist/track popularity, follower count,
release year, release type) over the user's listening corpus. The headline
metric is the area under the ROC curve (AUC); the engagement-magnitude
regression target that the lab opened on ($\text{play\_count}\times
\text{completion\_ratio}$, $\log1p$-transformed) is retained as a documented
scientific field but is no longer delivered. The behavioral co-listening
embedding and all playback/curation fields are deliberately excluded from the
feature set. Five Jun-2026 campaigns are folded in below; they converge on a
single finding --- categorical artist/genre/album identity read by
gradient-boosted trees beats every dense modality and every dense
distance/margin family --- and they crowned the project's first registered,
promotable champion, \texttt{xgboost-classifier-rotation}, at AUC
\num{0.7429} on the original DS-META corpus. Three Jun-12 campaigns then
re-ran the family question on a larger 2026 corpus (a full-64-d,
\num{23529}-row dataset): a four-family classifier bake-off in which
\texttt{xgboost-classifier} again wins (AUC \num{0.7204}), \texttt{catboost}
follows in-band, and a dense \texttt{pyramid-mlp} is last; a six-geometry
architecture scan that closes the dense-net family for AUC by \emph{refuting}
the budget hypothesis; and a topology bake-off of a new Rust/\texttt{burn}
\texttt{burn-deep-classifier} that, by feeding the identity one-hots
\emph{raw}, repairs the dense net's calibration completely (log-loss
$2.0\to 0.62$) and lifts its AUC $\sim$2 bands to \num{0.7038} --- yet still
lands $\sim$1--2 bands short of the tree tier, with learned entity embeddings
starved by the low identity cardinality. Section~\ref{sec:anns} consolidates why
artificial neural nets structurally underperform on this one-hot-heavy feature
matrix and how the representation-aware net resolves the diagnosis.
\end{abstract}

\section{Introduction}
This document is the living synthesis of the experiment campaigns; the
markdown reports in \texttt{experiments/} are the primary record. Each section
transcribes (never recomputes) the quantitative results of a campaign,
preserves its run and dataset identifiers verbatim for traceability, and adds
the rationale and interpretation that the terse reports omit.

\paragraph{Task and target.} The lab predicts \emph{taste fit} from a track's
intrinsic metadata over the user's listening corpus. It opened on an
engagement-magnitude regression target ---
$\text{play\_count}\times\text{completion\_ratio}$, trained against a
$\log1p$-transform and mapped back with $\text{expm1}$ --- but on 2026-06-08
the active target pivoted to a binary \emph{rotation} signal:
$\text{play\_count}\ge 2$, i.e.\ did the track earn a repeat listen
(Section~\ref{sec:rotation-pivot}). Rotation is the actionable input to the
Playlist Pathfinder's per-track desirability term, it is far more learnable
than engagement magnitude (trees reach $\text{AUC}\approx 0.73$ on rotation
versus \rsq{}$\approx 0.3$ on engagement), and with a base rate near
\SI{44.5}{\percent} it is naturally spread for path search. Engagement is
retained as a documented scientific field, but rotation is the delivered
target and every leaderboard below is on it.

\paragraph{Metrics.} Because the target is binary, the headline metric is the
area under the ROC curve (AUC, higher-better), the quantity the
classification leaderboard (Table~\ref{tab:lb-auc}) sorts on and the
best-models group auto-promotes by; it is reported alongside log-loss,
accuracy and the Brier score. Two distinct metric regimes nonetheless coexist
in the rotation campaigns. The campaigns opened with a \emph{regression} read
--- mean absolute error (MAE), RMSE and \rsq{} on the $[0,1]$ P(rotation)
scale --- and only later adopted the honest \emph{classification} read once
the framework gained a binary mode. MAE is retained on the historical
regression board (Table~\ref{tab:lb-mae}) for same-family band significance,
but it is documented there as \emph{degenerate} on a binary target: an
$\epsilon$-tube SVR can post a deceptively low MAE while being worse than
predicting the mean (Section~\ref{sec:svm}), an illusion the classification
metrics dissolve. The absolute-percentage metrics (MAPE/\medape{}) are
undefined on rotation, since half the labels are exactly zero. The historical
engagement board, where MAE \emph{is} the natural headline, is preserved in
\texttt{PROJECT-FACTS.md}; its numbers live on the $\log1p$ engagement scale
and are not comparable to anything below.

\paragraph{Features.} The feature space is a track's \emph{intrinsic metadata}
alone: genre, artist, artist/track popularity, follower count, release year,
credited-artist count and release type. The dense embedding modalities
(co-listening word2vec, content-text, acoustic, and the song-autoencoder
latent of Section~\ref{sec:ae-blend}) are \emph{behavioral or descriptive}
signals, not categorical identity, and the payload-based predictors over the
metadata one-hots are the canonical models here. Critically, the recipe is
\emph{leakage-free}: the playback fields that mechanically constitute or proxy
the target --- \texttt{play\_count}, \texttt{completion\_ratio},
\texttt{skip\_rate}, \texttt{shuffle\_rate}, \texttt{distinct\_sessions}, and
the demoted \texttt{engagement} field itself --- and the behaviorally
entangled curation fields are all excluded by default.

\paragraph{Two corpora.} The first three campaigns ran on the original
\emph{DS-META} corpus (\num{12256} rows, metadata-only). The two Jun-12
campaigns (Sections~\ref{sec:bakeoff-p64} and~\ref{sec:mlp-scan}) ran on a
larger 2026 corpus --- a \emph{full-64-d} dataset
(\dsid{ds-20260609-204419-p64-s42}, \num{23529} rows, 177 columns: 64 song-AE
PCA latent $+$ 7 numeric $+$ 106 identity one-hots). \textbf{The two corpora
are not comparable}: they differ in both size (\num{12256} versus \num{23529}
rows) and feature set (metadata-only versus $+$64 AE latent). AUC numbers on
the new corpus must therefore never be read against the \num{0.7429} DS-META
champion; each corpus carries its own leaderboard
(Tables~\ref{tab:lb-auc} and~\ref{tab:lb-auc-p64}) and nothing regressed when
the new-corpus xgboost lands at \num{0.7204}.

\section{Background: model families and metrics}
\label{sec:background}
This section gives the minimum theory for a reader who knows ML but not this
repo. It is keyed to \texttt{PROJECT-FACTS.md} and grows only when a new
predictor family or metric enters the reports.

\paragraph{The families in play.}
\begin{description}
  \item[Gradient-boosted trees (\texttt{xgboost}, \texttt{xgboost-classifier},
  \texttt{catboost-classifier}, \texttt{ngboost-classifier}).]
  An additive ensemble of shallow trees, each fit to the residual gradient of
  the previous stage. As a \emph{regressor} its output is a leaf average,
  bounded by the range of training targets; as a \emph{classifier} it emits a
  calibrated probability through a logistic link, the form that wins on
  rotation. Trees split on individual columns, so they read a track's sparse
  artist/genre/album identity one-hots \emph{losslessly} and capture non-linear
  interactions without manual feature crosses --- the reason they dominate this
  corpus. \texttt{catboost} is an ordered-boosting GBDT variant with symmetric
  (oblivious) trees and native categorical handling; on this corpus it is the
  legitimate second tree family, co-leader \emph{in-band} with xgboost.
  \texttt{ngboost} is a natural-gradient boosting variant that fits a full
  predictive \emph{distribution}; its probabilistic machinery, paid for with
  shallow base learners, costs it raw ranking power here. The cost of the GBDT
  family generally is a non-deterministic fit (subsampling and column sampling
  carry an internal RNG seed), so two runs at different seeds differ slightly.
  The registered champion \texttt{xgboost-classifier-rotation} is this family.
  \item[Logistic regression.] Linear log-odds passed through a sigmoid: a
  deterministic, $\ell_2$-regularized linear classifier. It is the natural
  \emph{linear floor} for a binary target --- the analogue of ridge in the
  regression world --- and serves as a cheap, reproducible control on the
  harness. Being a single linear boundary in a dense feature space, it blurs
  the sparse categorical identity that trees exploit, and it lands well below
  them on rotation.
  \item[Kernel SVM / SVC (\texttt{svc} with rbf / linear / poly / sigmoid
  kernels).] A maximum-margin classifier that operates on kernel-induced
  distances between points. As a \emph{regressor} (SVR) on a binary target the
  $\epsilon$-insensitive tube degenerates --- the source of the binary-MAE trap
  of Section~\ref{sec:svm}; as a proper probability \emph{classifier} (SVC) the
  divergence vanishes but every kernel still loses to the trees, because a
  dense-distance / large-margin model cannot read sparse identity one-hots the
  way an axis-aligned tree split does. The sigmoid kernel is the weakest and is
  the documented refutation echo.
  \item[Pyramid-MLP (\texttt{pyramid-mlp-classifier}).] A feed-forward
  artificial neural network with a tapering (``pyramid'') stack of hidden
  layers, ReLU activations and an Adam optimizer, fit on the StandardScaler-ed
  feature matrix. It is the project's only \emph{dense} (every-input-to-every-
  hidden-unit) learner on rotation. As Section~\ref{sec:anns} develops at
  length, a dense net is exactly the regime that one-hot-heavy taste data
  fights: it must learn dense weights over a mostly-zero input, and it lands
  last in the bake-off with a blown calibration that an architecture scan
  proves is structural, not a budget artifact.
  \item[Representation-aware net (\texttt{burn-deep-classifier}).] A
  full-Rust/\texttt{burn} feed-forward classifier built to repair the
  pyramid-MLP's pathologies. It splits the feature matrix into a
  \emph{standardized} continuous block (64 PCA $+$ 7 numeric) and a set of
  \emph{raw} (un-standardized) one-hot identity groups (artist, genre,
  album-type), so the scaler never turns rare categories into outlier spikes
  (Section~\ref{sec:anns-scaler}). A \texttt{topology} enum selects among a
  plain dense \texttt{mlp} over the concatenated inputs, learned per-group
  \texttt{embeddings}, and a \texttt{wide\_deep} split (a wide linear path over
  the one-hots beside a deep path over the continuous block). As
  Section~\ref{sec:anns-burn} reports, the raw-one-hot representation fixes the
  calibration completely and lifts AUC $\sim$2 bands over the pyramid, but the
  net still sits $\sim$1--2 bands below the tree tier, and the embedding
  variants gain nothing because the identity cardinalities here are too low to
  feed them.
  \item[Autoencoder-XGBoost blend (\texttt{autoencoder-xgboost}).] A weighted
  ensemble $\hat{y}=w\cdot\hat{y}_{\text{ae}}+(1-w)\cdot\hat{y}_{\text{xgb}}$
  that combines a regressor over a precomputed 64-dimensional
  song-autoencoder latent with a metadata-XGBoost model, the weight $w$ tuned on
  a train holdout. Its purpose was to test whether the dense AE latent carries
  rotation signal the metadata one-hots miss; it does not, and the tuner
  correctly drives $w\to 0$, collapsing the ensemble to metadata-only XGBoost
  (Section~\ref{sec:ae-blend}).
\end{description}

\paragraph{Why a binary target makes MAE degenerate.} For a continuous target,
mean absolute error is a faithful summary of how far predictions sit from the
truth. On a $\{0,1\}$ target it is not: a model can minimize MAE by hugging the
labels regardless of how it \emph{ranks} or \emph{calibrates} them. An SVR with
an $\epsilon{=}0.1$ tube snaps its outputs to a two-level step near
$\{0.1,0.9\}$, sitting only $\sim 0.1$ off every label --- a low MAE --- while
its ranking is worthless (negative \rsq{}). The metrics that matter for a
binary taste-fit signal are therefore the classification ones:
\begin{description}
  \item[AUC (area under the ROC curve).] The probability that a random positive
  is ranked above a random negative; \num{0.5} is chance, \num{1.0} is perfect.
  It is threshold-free and scale-free, depends only on the \emph{ordering} of
  the predicted probabilities, and is exactly the property the Pathfinder needs
  from its desirability term. It is the headline metric and the leaderboard
  sort key.
  \item[Log-loss.] The negative log-likelihood of the labels under the
  predicted probabilities; it punishes \emph{confident} mistakes severely and so
  rewards calibration, not just ranking.
  \item[Brier score.] The mean squared error between predicted probability and
  the $\{0,1\}$ label --- a calibration-and-refinement measure on the natural
  $[0,1]$ scale, lower-better.
\end{description}
A model that wins on AUC \emph{and} log-loss \emph{and} Brier, as the champion
does, is winning on ranking and calibration at once; the binary-MAE artifact
that flatters a degenerate SVR shows up on none of them. The same triad later
exposes the opposite pathology in a dense net (Section~\ref{sec:anns}): a model
can post a respectable AUC --- genuine ranking --- while its log-loss and Brier
are catastrophic, i.e.\ it ranks acceptably but is \emph{confidently wrong}.

\paragraph{A note on the historical log-target.} The lab's opening regression
target, engagement, is heavy-tailed and non-negative, so it was trained against
a $\log1p$-transform and mapped back with $\text{expm1}$. That transform pair
carries a hazard --- a model unbounded in log-space can produce a large
positive log prediction that $\text{expm1}$ amplifies into an enormous value,
the ``$\text{expm1}$ blow-up'' --- which is why bounded-by-construction families
(trees, rbf kernels) are preferred over unbounded ones (linear models, MLPs) on
that target. Engagement remains a documented field and its leaderboard lives in
\texttt{PROJECT-FACTS.md}, but it is no longer the delivered target, so the
blow-up hazard is now a historical footnote rather than a live design
constraint.

\section{The rotation pivot: from engagement magnitude to a binary taste-fit signal}
\label{sec:rotation-pivot}
On 2026-06-08 the active target pivoted away from engagement magnitude to a
binary \emph{rotation} signal, defined as $\text{play\_count} \ge 2$ --- did
the track earn a repeat listen? Rotation is the actionable input to the
Playlist Pathfinder's \texttt{W\_FIT} term, where a per-track $[0,1]$
desirability biases the A* path; a calibrated rotation probability is a
cleaner \texttt{fit\_norm} than a $\log1p$-then-$\text{expm1}$ engagement
estimate. Rotation also turns out to be far more learnable than engagement
magnitude (gradient-boosted trees reach $\text{AUC}\approx 0.73$ on rotation
versus \rsq{}$\approx 0.3$ on engagement) and is naturally spread for path
search, with a base rate near \SI{44.5}{\percent}. Engagement is retained as a
documented scientific field but is no longer the delivered target.

The pivot is decisive for how every number below must be read. Rotation lives
on the $[0,1]$ scale, so its MAE/\rsq{} are \emph{not} comparable to the
engagement leaderboard preserved in \texttt{PROJECT-FACTS.md}; and because half
the rotation labels are exactly zero, the absolute-percentage metrics
(MAPE/\medape{}) are undefined for those rows. Two distinct metric regimes
therefore coexist in the rotation campaigns: an early \emph{regression} read
(MAE, RMSE, \rsq{}, units P(rotation)) and the later, honest
\emph{classification} read (AUC, log-loss, accuracy, Brier). As the campaigns
below show, the choice of metric is not cosmetic --- under MAE a kernel SVR can
\emph{look} competitive while being worse than predicting the mean, an illusion
that the classification metrics dissolve.

\section{Rotation leaderboards}
\label{sec:leaderboard}
Three leaderboards. The first two are on the original DS-META corpus, one per
metric regime; the third is on the larger 2026 full-64-d corpus and is
\emph{not} comparable to the others (different corpus size and feature set ---
see the Introduction's ``Two corpora'' note). The DS-META AUC board carries the
project's first registered, promotable rotation champion; the new-corpus board
carries the new-corpus registered definition.

\begin{table}[ht]
  \centering
  \caption{Rotation, \emph{classification} regime, \textbf{DS-META corpus} ---
  ranked by AUC (higher-better). The champion is the only registered,
  auto-promoted definition. Datasets: DS-META
  \dsid{ds-20260609-141622-p1-s42} (115 cols, metadata-only, pca floor); the
  178-col row adds the 64 song-AE latent columns
  (\dsid{ds-20260608-204402-p64-s42}). AUC noise band $2\sigma=0.0126$.}
  \label{tab:lb-auc}
  \begin{tabular}{llcccc}
    \toprule
    Model & Dataset & AUC & log-loss & accuracy & Registered? \\
    \midrule
    \best{xgboost-classifier} (champion) & DS-META & \best{0.7429} & \best{0.5883} & 0.6846 & \best{yes (rank 1)} \\
    xgboost-classifier (178, +AE latent) & p64-s42 & 0.7350 & 0.6045 & 0.6818 & no \\
    svc (poly $d{=}3$, best non-tree)    & DS-META & 0.6855 & 0.6452 & 0.6450 & no \\
    logistic ($C{=}1$, linear floor)     & DS-META & 0.6762 & 0.6329 & 0.6491 & no \\
    svc (rbf)                            & DS-META & 0.6722 & 0.6469 & 0.6426 & no \\
    svc (linear)                         & DS-META & 0.6363 & 0.6535 & 0.6373 & no \\
    svc (sigmoid, worst)                 & DS-META & 0.6255 & 0.6673 & 0.6267 & no \\
    \bottomrule
  \end{tabular}
\end{table}

\begin{table}[ht]
  \centering
  \caption{Rotation, \emph{classification} regime, \textbf{new full-64-d
  corpus} --- ranked by AUC (higher-better). Dataset
  \dsid{ds-20260609-204419-p64-s42} (\num{23529} rows, 177 cols: 64 song-AE
  PCA $+$ 7 numeric $+$ 106 identity one-hots). \textbf{Not comparable to
  Table~\ref{tab:lb-auc}}: different corpus size and feature set; the only
  valid in-corpus reference is the \num{0.7204} xgboost point, which the
  bake-off reproduced byte-identically. The winner's default config is saved
  as the new-corpus definition. The single-seed sweep best (depth 8, lr 0.03,
  \num{0.7261}) is $+0.0057$ over the default --- inside the band --- so it is
  flagged for a 3-seed confirm, not promoted. The pyramid-mlp log-loss is its
  documented calibration ceiling (Section~\ref{sec:anns}); the
  \texttt{burn-deep-classifier} rows are the Rust/\texttt{burn}
  representation-aware net that repairs it (log-loss $\to\sim$0.62) yet still
  sits below the tree tier, so neither is registered
  (Section~\ref{sec:anns-burn}). AUC band $2\sigma\approx 0.0126$ (old-corpus
  estimate, approximate here).}
  \label{tab:lb-auc-p64}
  \begin{tabular}{llcccl}
    \toprule
    Model & max\_depth/lr & AUC & log-loss & accuracy & Note \\
    \midrule
    \best{xgboost-classifier} (winner, saved)      & 6 / 0.05  & \best{0.7204} & 0.6075 & 0.6585 & registered (new corpus) \\
    xgboost-classifier (sweep best, single-seed)   & 8 / 0.03  & 0.7261 & \best{0.6044} & \best{0.6651} & flagged for 3-seed confirm \\
    catboost-classifier (2nd tree, in-band)        & 6 / 0.05  & 0.7111 & 0.6132 & 0.6560 & $-0.0093$, inside band \\
    burn-deep-classifier (mlp, best tune)          & lr 0.0005 / [128,64] & 0.7038 & 0.6221 & 0.6521 & best ANN; calib.\ fixed \\
    burn-deep-classifier (mlp, Phase-1 winner)     & lr 0.001 / [128,64]  & 0.7001 & 0.6211 & 0.6470 & dense net; calib.\ fixed \\
    ngboost-classifier                             & ---       & 0.6917 & 0.6293 & 0.6392 & $\sim$2.3 bands back \\
    pyramid-mlp-classifier (arch-scan best)        & [512,256] & 0.6792 & 2.0584 & 0.6326 & dense net; calib.\ ceiling \\
    pyramid-mlp-classifier (bake-off, 300 iter)    & [256,128,64] & 0.6781 & 2.4085 & 0.6326 & dense net; calib.\ ceiling \\
    \bottomrule
  \end{tabular}
\end{table}

\begin{table}[ht]
  \centering
  \caption{Rotation, \emph{regression} regime, \textbf{DS-META corpus} ---
  ranked by MAE on the $[0,1]$ P(rotation) scale (lower-better). The
  metadata-only XGBoost leader is the \texttt{mae\_xgb} branch reproduced
  inside every blend run; it was never saved as a standalone definition. The
  SVM-poly row is listed for completeness and \emph{refuted}: its lower MAE is
  the $\epsilon$-tube artifact of Section~\ref{sec:svm}, its \rsq{} is
  negative. MAE noise band $2\sigma=0.0086$.}
  \label{tab:lb-mae}
  \begin{tabular}{lcccl}
    \toprule
    Model & MAE & RMSE & \rsq{} & Source \\
    \midrule
    \best{metadata-only XGBoost} (AE latent dropped) & \best{0.3879} & \best{0.4475} & \best{0.191} & ae-blend / svm \\
    xgboost on latent+metadata combined              & 0.3965 & --- & 0.166 & rotation pivot \\
    xgboost on AE-latent only (64 cols)              & 0.4098 & --- & 0.132 & rotation pivot \\
    svm (poly, $\epsilon$-tube artifact --- REFUTED) & 0.3814 & 0.5303 & $-0.137$ & svm \\
    \bottomrule
  \end{tabular}
\end{table}

\section{Campaigns}

\subsection{Autoencoder-latent + XGBoost blend: a dead modality, isolated}
\label{sec:ae-blend}
\emph{Rationale.} This was the first campaign on the pivoted rotation target.
The hypothesis was that a precomputed 64-dimensional song-autoencoder latent
(from the Qdrant collection \texttt{spotify\_tracks\_song\_ae}) carries a
track's \emph{identity} in a way the metadata one-hots do not, and that
blending an AE-latent regressor with a metadata-XGBoost model ---
$\hat{y}=w\cdot\hat{y}_{\text{ae}}+(1-w)\cdot\hat{y}_{\text{xgb}}$, with the
weight $w$ tuned on a train holdout --- could beat metadata-XGBoost alone. All
six arms ran on \dsid{ds-20260608-204402-p64-s42} (64 AE latent + 114 metadata
$=$ 178 cols, 12{,}256 rows, 9{,}805/2{,}451 at seed 42), varying only the AE
branch (kNN at $k\in\{10,25,50,100,200\}$, plus a ridge-on-latent arm).

\emph{Results.} All six runs (\runid{run-20260609-131545-e34e4} and siblings)
succeeded, and all returned byte-identical blended metrics --- MAE 0.38789,
\rsq{} 0.191, RMSE 0.4475 --- because the weight tuner drove $w\to 0$ in every
arm. The only axis that moved was the AE branch's standalone error
$mae_{\text{ae}}$, which grew monotonically with the neighborhood size, from
0.4259 at $k=10$ to 0.4702 at $k=200$ (the ridge arm landed at 0.4623, between
$k=100$ and $k=200$). Figure~\ref{fig:ae-knn} shows the curve against the three
horizontal references that matter: even the best AE read ($k=10$, 0.4259) never
reaches the metadata-only tree (0.3879), nor even the AE-latent-only
\emph{xgboost} on record (0.4098).

\begin{figure}[ht]
  \centering
  \includegraphics[width=0.8\textwidth]{ae_latent_knn}
  \caption{The song-AE latent carries no rotation signal the metadata
  one-hots do not already carry better. The kNN standalone error
  $mae_{\text{ae}}$ rises monotonically as the neighborhood widens (the
  latent's signal is strongest in tight local neighborhoods and washes toward
  the base rate), and even its best point sits above both the AE-latent-only
  tree and the metadata-only tree. The weight tuner correctly drives the blend
  weight to zero, so the ensemble collapses to metadata-only XGBoost.}
  \label{fig:ae-knn}
\end{figure}

\emph{Interpretation.} The collapse to $w=0$ is the correct answer, not a bug
to fix: given an AE branch strictly worse than the XGB branch on the holdout,
the optimal convex weight \emph{is} zero. The architecture's one real value was
that it \emph{isolated} the latent into a branch the tuner could discard,
recovering a cleaner result than feeding the latent into a unified XGBoost ---
which actually \emph{dilutes} the tree (MAE 0.3965 with the 64 latent columns
present versus 0.3879 with them dropped). The song-AE latent thus joins
co-listening, content-text and acoustic features as a dense modality that loses
to categorical identity one-hots. The one caveat at the time was honesty about
significance: the $0.3879$ versus $0.3965$ edge is only $0.0086$ MAE on a
$[0,1]$ target, and the rotation noise band was then unmeasured --- so the
campaign reported the number but asserted no significant win and promoted
nothing.
\emph{Source: \texttt{2026-06-09-ae-xgboost-blend.md}.}

\subsection{Kernel SVM on rotation: a refutation, and the binary-MAE trap}
\label{sec:svm}
\emph{Rationale.} With the latent ruled out, this campaign tested whether any
SVM kernel (rbf / linear / poly / sigmoid) produces a rotation predictor
competitive with the metadata-only XGBoost leader, and --- only if one passed
--- whether it added orthogonal signal in a grid-weighted blend. The metadata
matrix was reduced to the pca floor in a new dataset
(\dsid{ds-20260609-141622-p1-s42}, 115 cols), with two split-seed siblings
(seeds 17 and 101) built to measure the noise band, and the champion reproduced
as a control (\runid{run-20260609-141717-9a8ba-autoencoder-xgboost}, blend
weight 0). A secondary outcome was the campaign's central deliverable: the
\textbf{first measured rotation noise band}.

\emph{Results.} Across the three split seeds of the metadata-only XGBoost
control (s42 0.387891, s17 0.393072, s101 0.396417) the sample standard
deviation is $\sigma_{\text{split}}=0.00430$, giving a $2\sigma$ significance
band of \textbf{0.0086 MAE} on the $[0,1]$ scale. This retroactively settles
the previous campaign's open question: the $0.3879$ versus $0.3965$ edge is
$\Delta 0.0086 \approx 1{\cdot}\text{band}$ (likely better, not confident),
while the $0.3879$ versus $0.4098$ gap (AE-latent-only) is $\sim
2.5{\cdot}\text{band}$ and therefore real. As for the SVM arms, rbf, linear and
poly all posted MAE inside the control band (0.381--0.390), \emph{lower} than
the champion in the poly case (0.3814) --- yet every one of them is
worse-than-mean on \rsq{} ($-0.13$ to $-0.20$), with RMSE 0.53--0.54 against
the tree's 0.4475. Sigmoid SVR diverged outright: MAE 0.624, RMSE 1.531, \rsq{}
$-8.48$, with predictions reaching 24.7 on a $[0,1]$ target --- the classic
unbounded sigmoid-SVR blow-up. Figure~\ref{fig:svm} contrasts the two metric
views directly.

\begin{figure}[ht]
  \centering
  \includegraphics[width=0.95\textwidth]{svm_mae_vs_r2}
  \caption{On a binary target, low SVR MAE is an artifact, not skill. With
  $\epsilon=0.1$ on a $\{0,1\}$ target the SVR snaps its outputs to the
  $\epsilon$-tube edges (a hard two-level step, $p_{50}\approx 0.10$,
  $p_{90}\approx 0.90$), so each row sits only $\sim 0.1$ off its label and the
  \emph{MAE} (left) reads competitive. The honest \emph{$R^2$} (right) exposes
  it: every kernel is worse than predicting the mean. Sigmoid is off-scale at
  $R^2=-8.48$.}
  \label{fig:svm}
\end{figure}

\emph{Interpretation.} The low SVR MAE is a degeneracy of the metric on a
binary target. With an $\epsilon$-tube of width $0.1$ on $\{0,1\}$ labels, the
SVR collapses to a two-level step function pinned near the right classes ---
each row is about $0.1$ off its label, so the MAE is small, while calibration
and ranking are destroyed (negative \rsq{}). The metric that matters for the
Pathfinder \texttt{fit\_norm} (ranking and calibration, i.e.\ \rsq{}/RMSE/AUC)
says the SVM loses decisively. Worse, the two stage-2 blends
(\runid{run-20260609-143114-528ef-blend} and its lgbm twin) fitted weights of
exactly $[0.0, 1.0]$ --- the grid, optimizing the same degenerate validation
MAE, \emph{mis-selected pure SVM} and discarded the tree, the opposite of the
self-disabling seen in the blend campaign. The lesson is sharp: a grid-by-MAE
blend objective is itself unsafe on a binary target. Nothing was registered or
promoted; kernel SVM became the fifth dense-regime family to lose to one-hot
identity on this corpus.
\emph{Source: \texttt{2026-06-09-svm-kernel-rotation.md}.}

\subsection{The honest AUC bake-off and the first AUC champion}
\label{sec:auc}
\emph{Rationale.} The SVM campaign exposed that MAE is the wrong lens for a
binary target. The framework had meanwhile gained a classification mode
(\texttt{domain.toml} \texttt{task = binary}, primary metric AUC,
higher-better; the best-models group ranks by AUC and auto-promotes the top
run). This campaign re-ran the comparison on that honest footing: xgboost-classifier
against logistic and a four-kernel SVC scan, all on DS-META, plus a
178-column arm with the AE latent re-added to settle the latent-dilution
question on the AUC axis, and three control seeds to measure the AUC noise
band. The hypothesis --- a sixth test of the convergent finding --- was that
xgboost-classifier wins and every dense family loses.

\emph{Results.} All nine runs succeeded. The champion config
(\runid{run-20260609-172920-627d6-xgboost-classifier}, control hyperparameters,
seed 42) reached \textbf{AUC 0.7429} (log-loss 0.5883, accuracy 0.6846, Brier
0.2017) and was registered as \texttt{xgboost-classifier-rotation} and
auto-promoted to best-models rank 1 --- the project's first saved, AUC-scored,
promotable rotation champion. The first measured \emph{AUC} noise band comes
from the three control seeds (s42 0.742934, s17 0.730659, s101 0.734469, mean
0.736020): $\sigma_{\text{AUC}}=0.00628$, $2\sigma=\textbf{0.0126}$
(Figure~\ref{fig:auc-band}). This band is a different metric from the 0.0086
MAE band and does not replace it. Every dense family lost by many bands ---
svc-poly 0.6855, logistic 0.6762, svc-rbf 0.6722, svc-linear 0.6363,
svc-sigmoid 0.6255 --- a spread of 4.6 to 9.3 noise bands below the champion
(Figure~\ref{fig:auc-bakeoff}). The 178-column arm (AE latent re-added) scored
0.7350, $-0.0079$ below the metadata-only champion --- the same dilutive
direction as the regression result, but inside the AUC band, so
neutral-leaning-dilutive at a single seed.

\begin{figure}[ht]
  \centering
  \includegraphics[width=0.8\textwidth]{auc_bakeoff}
  \caption{The honest classification metric confirms the convergent finding a
  sixth time. xgboost-classifier reads the sparse artist/genre/album identity
  one-hots losslessly; the dense distance/margin families (logistic and all
  four SVC kernels) blur that identity in a dense feature space and land 4.6 to
  9.3 noise bands lower. The shaded strip is the $2\sigma=0.0126$ band below
  the champion; no challenger comes within a band of it.}
  \label{fig:auc-bakeoff}
\end{figure}

\begin{figure}[ht]
  \centering
  \includegraphics[width=0.7\textwidth]{auc_noise_band}
  \caption{The first measured rotation AUC noise band. Three data-split seeds
  of the registered champion config give $\sigma_{\text{AUC}}=0.0063$ and a
  $2\sigma$ significance band of $0.0126$. In relative terms this is actually
  tighter than the MAE band (0.85\% versus 1.10\% coefficient of variation),
  as expected for a rank metric on 2{,}451 test rows. The band only matters for
  the sub-band latent-dilution call; champion and challengers are separated by
  many bands.}
  \label{fig:auc-band}
\end{figure}

\emph{Interpretation.} Reframing rotation as classification did two things at
once. It dissolved the SVR illusion of Section~\ref{sec:svm} --- under AUC,
log-loss \emph{and} Brier, every SVC kernel is now visibly worse than the
trees, all in the same direction, so the $\epsilon$-tube artifact that made
SVR's MAE look competitive simply vanishes. And it confirmed that the metric
was only ever a problem for the \emph{fairness of the comparison}; the dense
family was always the loser. Notably svc-sigmoid does not collapse to chance
under classification (it holds AUC 0.6255 rather than blowing up): a proper
probability objective removes the unbounded divergence that wrecked the
SVR-sigmoid run. The convergent finding now stands on an honest metric: the
taste-fit signal lives in artist/genre/album \emph{identity}, captured by
sparse one-hots and exploited by trees, not in dense distance or margin models.
The remaining lever is identity \emph{resolution} --- raising the vocabulary
caps so the \texttt{artist=\_\_other\_\_} tail gets real columns --- not new
modalities or families. The next planned campaign is a champion hyperparameter
scan now that the family and metric are settled.
\emph{Source: \texttt{2026-06-09-rotation-classification.md}.}

\subsection{Four-family bake-off on the new full-64-d corpus, and an xgboost depth/lr sweep}
\label{sec:bakeoff-p64}
\emph{Rationale.} The first three campaigns ran on the original DS-META corpus.
This campaign re-asked the family question on the larger 2026 corpus --- ``all
of the 64-d embedding'', the full 177-column matrix
(\dsid{ds-20260609-204419-p64-s42}, \num{23529} rows, 64 song-AE PCA $+$ 7
numeric $+$ 106 identity one-hots) --- by putting four classifier families head
to head at registry defaults, seed 42: \texttt{xgboost-classifier},
\texttt{catboost-classifier}, \texttt{ngboost-classifier} and a dense
\texttt{pyramid-mlp-classifier}. The latter three are \emph{new families on
this project}. The winner then earned a \texttt{max\_depth}$\times$
\texttt{learning\_rate} sweep. The only valid in-corpus reference is a
pre-existing run on this dataset (\runid{run-20260609-204726-439d8-xgboost-classifier},
AUC \num{0.7204}); nothing here may be read against the DS-META \num{0.7429}
champion (different corpus and feature set).

\emph{Results.} All four bake-off arms succeeded (Table~\ref{tab:bakeoff-p64};
Figure~\ref{fig:bakeoff-p64}). \texttt{xgboost-classifier}
(\runid{run-20260612-115909-840c0-xgboost-classifier}) won at AUC \num{0.7204},
reproducing the in-corpus reference byte-identically (\num{0.7204383403}) ---
the determinism gate passed. \texttt{catboost-classifier}
(\runid{run-20260612-115909-c7ef6-catboost-classifier}) followed at \num{0.7111},
$-0.0093$ behind --- \emph{inside} the $\sim$0.0126 band, so co-leader and the
legitimate second tree family, not a loss. \texttt{ngboost-classifier}
(\runid{run-20260612-115909-d725b-ngboost-classifier}) landed at \num{0.6917}
($\sim$2.3 bands back), and the dense \texttt{pyramid-mlp-classifier}
(\runid{run-20260612-115909-ea698-pyramid-mlp-classifier}) was last at
\num{0.6781} --- with a log-loss of \num{2.4085} and Brier \num{0.3251} far
above the trees' $\sim$0.60--0.63 / $\sim$0.21, the diagnostic dissected in
Section~\ref{sec:anns}.

\begin{table}[ht]
  \centering
  \caption{Phase 1 --- the four-family bake-off on the new full-64-d corpus
  (\dsid{ds-20260609-204419-p64-s42}, seed 42, registry defaults). All four
  succeeded. AUC higher-better; log-loss and Brier lower-better. Winner bolded;
  best cell per metric bolded.}
  \label{tab:bakeoff-p64}
  \begin{tabular}{llcccc}
    \toprule
    Arm / run id & Predictor & AUC & log-loss & accuracy & Brier \\
    \midrule
    \best{C0} (\runid{run-20260612-115909-840c0-xgboost-classifier}) & \best{xgboost-classifier}   & \best{0.7204} & \best{0.6075} & \best{0.6585} & \best{0.2109} \\
    A1 (\runid{run-20260612-115909-c7ef6-catboost-classifier})       & catboost-classifier         & 0.7111 & 0.6132 & 0.6560 & 0.2132 \\
    A2 (\runid{run-20260612-115909-d725b-ngboost-classifier})        & ngboost-classifier          & 0.6917 & 0.6293 & 0.6392 & 0.2193 \\
    A3 (\runid{run-20260612-115909-ea698-pyramid-mlp-classifier})    & pyramid-mlp-classifier      & 0.6781 & 2.4085 & 0.6326 & 0.3251 \\
    \bottomrule
  \end{tabular}
\end{table}

\begin{figure}[ht]
  \centering
  \includegraphics[width=0.8\textwidth]{bakeoff_p64}
  \caption{The family ordering survives the corpus change. On the larger 2026
  full-64-d matrix the two gradient-boosted-tree families lead
  (\texttt{xgboost} winning, \texttt{catboost} co-leader inside the shaded
  $2\sigma$ band below the winner); \texttt{ngboost} sits $\sim$2.3 bands back;
  the dense \texttt{pyramid-mlp} is last, the same convergent finding seen on
  DS-META in Figure~\ref{fig:auc-bakeoff}.}
  \label{fig:bakeoff-p64}
\end{figure}

The winner earned a Phase-2 sweep over \texttt{max\_depth}$\in\{4,6,8\}\times$
\texttt{learning\_rate}$\in\{0.03,0.05,0.1\}$ (Table-form data in the source
report; Figure~\ref{fig:xgb-heatmap}). Depth is the dominant axis: every
depth-4 cell is at the bottom (0.7074--0.7119), and the depth-8 / lr-0.03 corner
tops the grid at \num{0.7261} (\runid{run-20260612-120207-2dd3d-xgboost-classifier},
log-loss \num{0.6044}). High learning rate (0.1) consistently hurts both AUC and
calibration. But the best-vs-default gap (+0.0057) is \emph{inside} the band, so
the default depth-6/lr-0.05 config was saved as the new-corpus definition and the
depth-8/lr-0.03 cell was flagged for a 3-seed confirmation rather than promoted.

\begin{figure}[ht]
  \centering
  \includegraphics[width=0.62\textwidth]{xgb_depth_lr_heatmap}
  \caption{The xgboost \texttt{max\_depth}$\times$\texttt{learning\_rate} sweep
  on the new corpus. Depth is the lever (the AUC climbs monotonically up the
  depth axis at the slow learning rates), and a fast learning rate hurts. The
  best cell (depth 8, lr 0.03, boxed) reaches \num{0.7261} but only $+0.0057$
  over the default --- inside the $2\sigma$ band --- so the default config
  remains saved and the best cell awaits a 3-seed confirm.}
  \label{fig:xgb-heatmap}
\end{figure}

\emph{Interpretation.} The convergent finding survives the corpus change intact:
on a different, larger corpus with the AE latent present, the two tree families
again lead and the dense net is again last. \texttt{catboost} earns its place as
a legitimate second tree family (co-leader in-band) --- a tree-diversity /
second-opinion member, not an AUC gain over xgboost. \texttt{ngboost} pays for
its distributional machinery: shallow base learners (\texttt{base\_max\_depth} 3)
under-fit the 177-column one-hot-heavy matrix. The pyramid-MLP's last-place AUC
\emph{plus} its blown log-loss are the dense-net signature; the bake-off read the
log-loss as a 300-iteration under-training artifact and prescribed more budget ---
a hypothesis the next campaign tests and refutes.
\emph{Source: \texttt{2026-06-12-classifier-bakeoff-p64.md}.}

\subsection{Pyramid-MLP architecture scan: refuting the budget hypothesis}
\label{sec:mlp-scan}
\emph{Rationale.} The bake-off flagged the pyramid-MLP's last place (AUC
\num{0.6781}, log-loss \num{2.4085}) as a 300-iteration-cap under-training
artifact, prescribing ``raise \texttt{max\_iter} so the net converges.'' This
follow-up tests that mechanism directly: six pyramid geometries, all with
\texttt{max\_iter} raised to \textbf{800}, alpha \num{1e-4}, learning rate
\num{1e-3}, seed 42, on the same dataset --- so the comparison reflects
\emph{architecture}, not budget. Does any geometry, given the budget to
converge, lift the dense family to a competitive AUC?

\emph{Results.} Refuted on every count (Table~\ref{tab:mlp-scan};
Figures~\ref{fig:mlp-arch} and~\ref{fig:mlp-calib}). \textbf{No arm hit the 800
cap}: every geometry self-stopped at \textbf{73--100 iterations} via sklearn's
\texttt{n\_iter\_no\_change} loss-plateau tolerance --- far below even the old
300 cap, so the budget was never the binding constraint. The decisive tell: the
re-baseline arm M0 ([256,128,64], the bake-off geometry now at \texttt{max\_iter}
800) reproduced the under-trained run's AUC to all digits ---
\num{0.6780926} versus \num{0.6780926} --- because it ran the identical 73
iterations. Raising the cap from 300 to 800 changed nothing. The best converged
geometry was M1 [512,256] at AUC \num{0.6792} (\runid{run-20260612-123155-0ace5}),
a statistical tie with the under-trained default ($+0.0011 \approx$ zero) and
$\sim$3.3 bands below xgboost. The whole family lands in a 0.0080-wide AUC strip
(0.6712--0.6792), entirely inside one noise band --- the geometry ranking is not
resolvable at a single seed --- with the log-loss (2.06--2.78) and Brier
(0.32--0.34) blown at the same level as the under-trained default across every
arm. No definition was saved (below the 0.683 save bar).

\begin{table}[ht]
  \centering
  \caption{The six-geometry pyramid-MLP architecture scan
  (\dsid{ds-20260609-204419-p64-s42}, \texttt{max\_iter} 800, alpha \num{1e-4},
  lr \num{1e-3}, seed 42), sorted by AUC. All six succeeded; \emph{none} hit
  the 800-iteration cap --- each self-stopped at the listed iteration count.
  Best cell per metric bolded.}
  \label{tab:mlp-scan}
  \begin{tabular}{llcccc}
    \toprule
    Arm & Hidden layers & AUC & log-loss & Brier & Iters (of 800) \\
    \midrule
    \best{M1} wide-shallow & [512, 256]            & \best{0.6792} & \best{2.0584} & \best{0.3216} & 75 \\
    M0 re-baseline         & [256, 128, 64]        & 0.6781 & 2.4085 & 0.3251 & 73 \\
    M4 gentle taper        & [384, 307, 246]       & 0.6777 & 2.7764 & 0.3292 & 87 \\
    M2 deep pyramid        & [256, 128, 64, 32]    & 0.6737 & 2.7631 & 0.3303 & 100 \\
    M3 steep taper         & [512, 169, 56]        & 0.6736 & 2.7031 & 0.3365 & 89 \\
    M5 wide-deep           & [512, 307, 184, 111]  & 0.6712 & 2.6457 & 0.3348 & 84 \\
    \bottomrule
  \end{tabular}
\end{table}

\begin{figure}[ht]
  \centering
  \includegraphics[width=0.95\textwidth]{mlp_arch_scan}
  \caption{Both knobs of the dense family are exhausted. \emph{Left:} the six
  geometries land inside a single $2\sigma$ band around the under-trained
  default --- architecture is not the binding constraint. \emph{Right:} every
  geometry self-stops at 73--100 iterations, well below even the old 300 cap
  (let alone the raised 800), so the training \emph{budget} was never the
  binding constraint either. The optimizer reaches a loss plateau it cannot
  escape at this learning rate and halts. The binding constraint is the input
  \emph{representation}.}
  \label{fig:mlp-arch}
\end{figure}

\emph{Interpretation.} The refutation is clean and twofold: neither the training
budget (every arm self-stops far below the cap; M0 reproduces the under-trained
AUC to all digits) nor the architecture (all six geometries inside one band)
moves the dense family. The high log-loss is therefore not a truncated training
trajectory but a genuine \emph{calibration ceiling} at these hyperparameters ---
the optimizer's chosen plateau. With the dense family's own knobs exhausted, the
binding constraint must lie elsewhere: in the input \emph{representation}, the
subject of Section~\ref{sec:anns}. The directional within-family trend ---
wider-and-shallower mildly favored (M1 top), depth/capacity mildly harmful (M5
wide-deep last) --- is consistent with a dense net struggling to route signal
through 106 sparse one-hots: extra layers add capacity to memorize and blur,
not to separate. The family is closed for AUC on this feature set; do not
re-scan geometry or budget.
\emph{Source: \texttt{2026-06-12-pyramid-mlp-arch-scan.md}.}

\section{Why artificial neural nets underperform on one-hot-heavy taste data}
\label{sec:anns}
The dense \texttt{pyramid-mlp} is the only artificial neural net in the
rotation campaigns, and across two Jun-12 campaigns it lands last with a
distinctive failure pattern. This section consolidates, rigorously, \emph{why}
--- grounding every claim in the campaign evidence. The short version is the
project's documented \emph{convergent finding}: on this feature matrix the
binding constraint is the input \emph{representation}, and a dense net is the
one model class structurally mismatched to it. We build the argument in five
parts and close with the motivated next step.

\subsection{The feature matrix is mostly sparse identity}
\label{sec:anns-matrix}
The full-64-d matrix has 177 columns, and their \emph{composition} is the whole
story. \textbf{106 of the 177 columns are sparse identity one-hots}: artist
($\times 61$), \texttt{genre\_primary} ($\times 41$) and \texttt{album\_type}
($\times 4$). Only 71 columns are continuous --- the 64 song-AE PCA latents and
7 numeric metadata fields (popularities, follower count, release year,
credited-artist count). For any given track exactly one artist column, one genre
column and one album-type column is hot; the rest of the 106 are zero. The
matrix is, by construction, dominated by a mostly-zero block of mutually
exclusive indicators.

A gradient-boosted tree reads this block \emph{losslessly and cheaply}: a single
axis-aligned split on column \texttt{artist=X} partitions the data on exactly
that identity, with no interaction among the one-hots forced and no cost paid for
the columns that are zero. Capturing ``tracks by artist $X$ in genre $Y$ rotate''
costs the tree two nested splits. A dense net cannot do this. It must learn a
\emph{dense} weight from every one of the 106 mostly-zero inputs into every
hidden unit, then back out a non-linear combination --- spending capacity and
gradient signal modeling the structure of a block that the tree handles with one
comparison. This is precisely the regime in which trees structurally dominate
dense nets, and it is the project's \emph{convergent finding}, now observed
across four feature modalities and five model families: the taste-fit signal is
carried by artist/genre identity, captured by one-hots and exploited by trees,
not by descriptive attributes or distance/margin/dense models. The pyramid-MLP
is simply the dense family's turn to lose to it.

\subsection{Standardization turns rare categories into outlier spikes}
\label{sec:anns-scaler}
The dense net is fed a StandardScaler-ed matrix, and standardization interacts
pathologically with sparse one-hots. A one-hot column for a category present in
a fraction $p$ of rows has mean $p$ and standard deviation $\sqrt{p(1-p)}$.
Standardizing it maps the ``on'' rows to
\[
  z_{\text{on}} = \frac{1-p}{\sqrt{p(1-p)}} = +\sqrt{\frac{1-p}{p}},
  \qquad
  z_{\text{off}} = \frac{0-p}{\sqrt{p(1-p)}} = -\sqrt{\frac{p}{1-p}}.
\]
For a \emph{rare} category --- the long tail of artists and genres, where $p$ is
small --- the ``on'' value $+\sqrt{(1-p)/p}$ becomes a large positive spike (for
$p=0.01$ it is $\approx +9.9$), while the ``off'' value sits at a tiny negative
$-\sqrt{p/(1-p)}\approx -0.1$. Every rare identity therefore arrives at the net
as an occasional, enormous outlier on its own axis. ReLU units with Adam at
learning rate \num{1e-3} swing hard on these spikes --- a single rare-category
activation drives a large pre-activation that the optimizer chases --- and the
softmax/sigmoid head saturates, emitting near-0 or near-1 probabilities on thin
evidence. Trees are immune: a split threshold on a one-hot is invariant to any
monotone rescaling, so the scaler's spike does nothing to them. This is a direct
mechanistic contributor to the blown calibration documented next.

\subsection{The diagnostic signature: genuine ranking, broken calibration}
\label{sec:anns-signature}
The dense net's failure has a precise, reproducible signature in the metrics ---
and it is \emph{not} the failure of a weak model. The pyramid-MLP posts AUC
$\approx 0.68$ (0.6781 in the bake-off, 0.6792 best in the arch scan): genuine
ranking signal, well above the 0.5 chance line, only $\sim$3 bands below the
trees. Yet \emph{at the same time} its log-loss is $\approx 2.0$--$2.8$ and its
Brier $\approx 0.32$--$0.34$, against the trees' log-loss $\approx 0.60$--$0.63$
and Brier $\approx 0.21$ (Figure~\ref{fig:mlp-calib}). A constant base-rate
predictor (always outputting $p\approx 0.445$) scores a log-loss of only
$\approx 0.69$. \textbf{The net's log-loss of 2.0--2.8 is thus far worse than the
random baseline, while xgboost's 0.61 is well below it.}

\begin{figure}[ht]
  \centering
  \includegraphics[width=0.95\textwidth]{mlp_calibration_signature}
  \caption{The dense-net diagnostic: the net is confidently wrong, not merely
  weak. \emph{Left:} on AUC the pyramid-MLP genuinely ranks ($\approx 0.68$,
  well above chance, $\sim$3 bands below the trees) --- there is real signal in
  its ordering. \emph{Right:} but its log-loss ($\approx 2.0$--$2.8$) is far
  \emph{worse} than the random base-rate baseline ($\approx 0.69$, dashed),
  while the trees sit well below it ($\approx 0.61$). A model that ranks
  acceptably yet scores worse-than-random log-loss is emitting confident,
  miscalibrated probabilities.}
  \label{fig:mlp-calib}
\end{figure}

This combination --- acceptable ranking, worse-than-random log-loss --- means the
net is \emph{confidently wrong}: it pushes many predictions to the saturated
extremes (near 0 or near 1) on the wrong side of the label, and log-loss punishes
each such confident error severely. It is miscalibrated, not weak: the ordering
is there, but the probabilities attached to that ordering are unusable as a
\texttt{fit\_norm}. The mechanism is exactly the scaler-spike saturation of
Section~\ref{sec:anns-scaler}. Contrast the trees, whose leaf-frequency outputs
are calibrated by construction: a leaf reporting probability $q$ is the empirical
rotation frequency of the training rows that fall in it, so the tree's
probabilities are honest by design --- the reason it wins on log-loss and Brier as
well as AUC.

\subsection{Optimization versus representation: the binding constraint}
\label{sec:anns-binding}
A natural objection is that the net is merely under-optimized --- too few
iterations, the wrong geometry. The architecture scan (Section~\ref{sec:mlp-scan})
forecloses both. On the \emph{budget} axis, every one of the six geometries
self-stops at 73--100 iterations, well below even the old 300 cap (let alone the
raised 800): the optimizer reaches a loss plateau it cannot escape, so more budget
buys nothing --- the re-baseline arm M0 reproduced the under-trained AUC to all
digits because it ran the identical 73 iterations. On the \emph{architecture}
axis, all six geometries land inside a single noise band, and the best converged
geometry reproduces the under-trained AUC to four digits (0.6792 versus 0.6781).
Both knobs internal to the dense family are therefore exhausted. By elimination,
the binding constraint is neither budget nor architecture but the input
\emph{representation} --- the raw one-hots of Section~\ref{sec:anns-matrix}, which
the net cannot consume the way a tree does. This is the diagnostic that points at
representation-aware topologies (entity embeddings, wide-and-deep) as the only
untested lever, since those are designed precisely to replace raw one-hots with a
learned dense representation.

\subsection{PCA has already pre-empted the net's nonlinear job}
\label{sec:anns-pca}
There is a second reason the dense net has little to do. The 64 song-AE PCA
latents are already a \emph{linear, decorrelated} compression of a
1536-dimensional embedding: PCA has extracted the dominant linear structure and
handed the net orthogonal components with most of the easy linear variance
already separated. Whatever nonlinear interaction a hidden layer might discover
in that block has largely been pre-empted; little nonlinear structure remains for
the net's depth to exploit. The evidence is direct: a plain logistic regression
reaches AUC 0.676 on this corpus and the converged MLP reaches 0.679 --- a dead
heat. A multi-layer ReLU network buys essentially nothing over a single linear
boundary, which is exactly what one expects when the continuous block is already
linearly pre-conditioned and the categorical block is the part that carries
signal the net cannot read.

\subsection{The open question and the motivated next step}
\label{sec:anns-next}
The diagnosis is now sharp: the dense net's ceiling is set by its inability to
consume raw one-hot identity, not by budget, architecture, or the continuous
block. That isolates a single, testable open question --- \emph{can a
representation-aware topology close the gap to the trees?} The two ingredients
that target the diagnosis directly are (i) \textbf{entity embeddings} for the
identity groups (artist, genre, album-type), which replace each sparse one-hot
block with a learned low-dimensional dense vector --- the standard fix for
high-cardinality categoricals in nets, removing both the dense-weights-over-zeros
problem of Section~\ref{sec:anns-matrix} and the scaler-spike problem of
Section~\ref{sec:anns-scaler}; and (ii) a \textbf{wide-and-deep} split --- a wide
linear path over the one-hots (so identity is read directly, as the tree reads
it) alongside a deep path over the continuous block. A Rust/\texttt{burn}
implementation of such a representation-aware network --- the
\texttt{burn-deep-classifier} --- was built precisely to answer this question.
Section~\ref{sec:anns-burn} reports its results, and they resolve the diagnosis
cleanly: the calibration pathology was indeed a representation artifact and is
fixed, but the AUC gap to the trees only \emph{narrows}, it does not close.

\emph{Sources: \texttt{2026-06-12-classifier-bakeoff-p64.md},
\texttt{2026-06-12-pyramid-mlp-arch-scan.md}.}

\subsection{Resolved: the burn-deep-classifier fixes calibration but not the AUC gap}
\label{sec:anns-burn}
\emph{Rationale.} The preceding subsections left a single sharp open question:
the dense net's ceiling is set by its inability to consume raw one-hot identity,
so \emph{can a representation-aware topology close the gap to the trees?} The
\texttt{burn-deep-classifier} --- a full-Rust/\texttt{burn} network --- was built
to answer it. Its design targets the two diagnosed pathologies directly: it keeps
the identity one-hots \emph{raw} (standardizing only the 64 PCA $+$ 7 numeric
continuous block), so the scaler-spike mechanism of Section~\ref{sec:anns-scaler}
never fires, and it exposes a \texttt{topology} enum over a plain dense
\texttt{mlp}, learned entity \texttt{embeddings}, and a \texttt{wide\_deep} split.
Phase 1 compared the three topologies (epochs 250, patience 20, early-stop on
validation log-loss, hidden $[128,64]$, embed\_dim 8, lr \num{1e-3}, batch 256);
Phase 2 tuned the winner. All nine runs ran on the same full-64-d dataset
(\dsid{ds-20260609-204419-p64-s42}) and succeeded cleanly.

\emph{Results.} The headline is the calibration recovery
(Figure~\ref{fig:burn-topology}). Every topology lands at log-loss
\num{0.621}--\num{0.628} and Brier \num{0.217}--\num{0.220} --- squarely in the
trees' $\sim$0.60--0.63 / $\sim$0.21 regime --- against the sklearn pyramid's
log-loss \num{2.0}--\num{2.4} and Brier \num{0.33}. \textbf{The calibration
pathology is fixed completely.} On the topology axis (Table~\ref{tab:burn-topo}),
the plain \texttt{mlp} wins at AUC \num{0.7001}, with \texttt{embeddings}
(\num{0.6934}) and \texttt{wide\_deep} (\num{0.6918}) both inside the
$2\sigma=0.0126$ band of it and of each other --- \texttt{embeddings} edges
\texttt{wide\_deep} by a sub-band $+0.0016$. The Phase-2 \texttt{mlp} tune
(Table~\ref{tab:burn-tune}) tops out at AUC \num{0.7038} (lr \num{5e-4},
$[128,64]$): the whole 6-cell grid is a \num{0.0081}-wide strip
(\num{0.6957}--\num{0.7038}), every $[128,64]$ cell beats every $[256,128]$ cell,
and the best config is only $+0.0037$ (inside the band) over the Phase-1 winner,
so the Phase-1 \texttt{mlp} config stands. At \num{0.7038} the net sits
$-0.0073$ below \texttt{catboost} ($\sim$0.6 band) and $-0.0166$ below
\texttt{xgboost} ($\sim$1.3 bands); below the $\sim$0.71 save bar, so no
definition was persisted.

\begin{table}[ht]
  \centering
  \caption{Phase 1 --- the \texttt{burn-deep-classifier} topology bake-off
  (\dsid{ds-20260609-204419-p64-s42}, held params: epochs 250, patience 20,
  hidden $[128,64]$, embed\_dim 8, dropout 0.1, lr \num{1e-3}, batch 256, relu;
  the only axis is \texttt{topology}). All three succeeded. AUC higher-better;
  log-loss and Brier lower-better. Winner bolded; best cell per metric bolded.}
  \label{tab:burn-topo}
  \begin{tabular}{llcccc}
    \toprule
    Arm / run id & Topology & AUC & log-loss & accuracy & Brier \\
    \midrule
    \best{T-mlp} (\runid{run-20260612-140116-41fc9-burn-deep-classifier}) & \best{mlp}        & \best{0.7001} & \best{0.6211} & \best{0.6470} & \best{0.2169} \\
    T-emb (\runid{run-20260612-140116-0eda4-burn-deep-classifier})        & embeddings        & 0.6934 & 0.6251 & 0.6436 & 0.2187 \\
    T-wd (\runid{run-20260612-140116-e85be-burn-deep-classifier})         & wide\_deep        & 0.6918 & 0.6281 & 0.6407 & 0.2198 \\
    \bottomrule
  \end{tabular}
\end{table}

\begin{table}[ht]
  \centering
  \caption{Phase 2 --- tuning the winning \texttt{mlp} topology, a
  \texttt{lr}$\times$\texttt{hidden} grid (other params held). All six
  succeeded, sorted by AUC; winner bolded, best cell per metric bolded. The
  whole grid is inside one noise band, every $[128,64]$ cell beats every
  $[256,128]$ cell, and the best gain over the Phase-1 winner ($+0.0037$) is
  sub-band --- so the Phase-1 config is the config of record.}
  \label{tab:burn-tune}
  \begin{tabular}{lllcccc}
    \toprule
    Run id & lr & hidden & AUC & log-loss & accuracy & Brier \\
    \midrule
    \best{\runid{run-20260612-140211-6575a-burn-deep-classifier}} & \best{0.0005} & \best{[128,64]}  & \best{0.7038} & 0.6221 & 0.6521 & 0.2172 \\
    \runid{run-20260612-140211-21ad5-burn-deep-classifier}        & 0.002         & [128,64]         & 0.6997 & 0.6227 & \best{0.6532} & 0.2174 \\
    \runid{run-20260612-140211-18c9e-burn-deep-classifier}        & 0.001         & [128,64]         & 0.6996 & \best{0.6213} & 0.6447 & \best{0.2170} \\
    \runid{run-20260612-140211-fecb0-burn-deep-classifier}        & 0.0005        & [256,128]        & 0.6966 & 0.6232 & 0.6424 & 0.2179 \\
    \runid{run-20260612-140211-4c359-burn-deep-classifier}        & 0.002         & [256,128]        & 0.6963 & 0.6245 & 0.6460 & 0.2181 \\
    \runid{run-20260612-140211-0f4e1-burn-deep-classifier}        & 0.001         & [256,128]        & 0.6957 & 0.6252 & 0.6424 & 0.2187 \\
    \bottomrule
  \end{tabular}
\end{table}

\begin{figure}[ht]
  \centering
  \includegraphics[width=0.95\textwidth]{burn_topology}
  \caption{The representation-aware net resolves the open question of
  Section~\ref{sec:anns-next}. \emph{Left:} the three topologies are inside one
  $2\sigma$ band of each other; the plain \texttt{mlp} wins and the entity
  \texttt{embeddings} barely edge \texttt{wide\_deep}, because at these identity
  cardinalities (artist 61, genre 41, album-type 4) learned embeddings are
  starved --- they pay off at hundreds-to-thousands of categories, not dozens.
  \emph{Right:} the headline. Keeping the one-hots \emph{raw} drops the
  log-loss from the sklearn pyramid's $\approx 2.0$ (confidently miscalibrated,
  worse than the random baseline) to $\approx 0.62$, into the tree regime --- a
  clean confirmation that the calibration pathology was a representation
  artifact, not an intrinsic property of nets on this data.}
  \label{fig:burn-topology}
\end{figure}

\emph{Interpretation.} The result resolves Section~\ref{sec:anns-next}'s open
question in two parts, and the two answers point in opposite directions. First,
the \emph{calibration} verdict is overturned: the pyramid-MLP's worse-than-random
log-loss of Section~\ref{sec:anns-signature} was a \emph{representation} artifact
--- raw one-hots crushed into the standardizer --- not an intrinsic property of
dense nets on taste data, and a net that keeps the one-hots raw is calibrated in
the trees' regime. The earlier ``broken calibration is a fixed ceiling'' reading
holds only for the sklearn-pyramid representation. AUC also recovered: the burn
net's \num{0.7038} is $+0.025$ over the best converged pyramid (\num{0.6792}),
$\sim$2 bands of genuine ranking-power gain, so the prior conclusion that ``the
MLP family is closed for AUC at $\approx 0.68$'' is revised --- a raw-one-hot
Rust net reaches $\sim$0.70. Second, and against the optimistic reading, the
representation fix does \emph{not} reach the tree tier: at \num{0.7038} the net is
still $\sim$1 band below \texttt{catboost} and $\sim$1.3 below \texttt{xgboost}.
The gap narrowed dramatically (from the pyramid's $\sim$3.3 bands to $\sim$1) but
did not close, so the convergent finding survives intact --- trees still read the
sparse identity one-hots a hair more effectively than even a raw-one-hot net. The
topology comparison sharpens this: the plain \texttt{mlp} beating both the
embedding and wide-deep variants is the \emph{same} ``the identity is already in
the one-hots'' lesson the tree campaigns taught, now reproduced inside the ANN
family --- learned per-group embeddings have little to add over 61 artist and 41
genre one-hots, and the standing lever for the net (as for the trees) is identity
\emph{resolution}: raise the vocabulary caps and the embedding topology should
finally start to pay off. No definition was saved (below the $\sim$0.71 bar); the
net is a strong, now-honest dense baseline and a natural ensemble member for its
decorrelated inductive bias, not a champion.
\emph{Source: \texttt{2026-06-12-burn-deep-topology.md}.}

\section{Conclusions}
\label{sec:conclusions}
Six Jun-2026 campaigns on the pivoted rotation target converge on a single
story. (1) Across two feature campaigns and three model-family scans (two
corpora), every dense modality and every dense distance/margin/dense-net family
loses to categorical identity one-hots read by gradient-boosted trees --- the
song-AE latent dilutes (Section~\ref{sec:ae-blend}), kernel SVM is
worse-than-mean (Section~\ref{sec:svm}), and on the new full-64-d corpus the two
tree families lead while the dense pyramid-MLP is last
(Section~\ref{sec:bakeoff-p64}). (2) The choice of metric is load-bearing on a
binary target: MAE is degenerate (a two-level $\epsilon$-tube predictor posts a
deceptively low MAE while being worse than the mean), and a grid-by-MAE blend
even mis-selects such a predictor; the honest classification metrics (AUC,
log-loss, Brier) restore a fair comparison (Section~\ref{sec:auc}). (3) The
campaigns delivered two measured noise bands --- $2\sigma=0.0086$ MAE and
$2\sigma=0.0126$ AUC --- and the project's first registered, promotable rotation
champion, \texttt{xgboost-classifier-rotation} at AUC 0.7429 on DS-META, plus a
saved new-corpus xgboost definition at AUC 0.7204. (4) The dense neural net was
diagnosed in full (Section~\ref{sec:anns}): an architecture scan refuted the
budget hypothesis (every geometry self-stops at 73--100 iterations and the
converged net reproduces the under-trained AUC to all digits), so the binding
constraint is the raw one-hot \emph{representation}, not optimization --- and the
net's worse-than-random log-loss beside a genuine AUC marks it as confidently
miscalibrated, not merely weak. (5) That diagnosis was then tested directly: the
Rust/\texttt{burn} \texttt{burn-deep-classifier}, which feeds the identity
one-hots \emph{raw}, repaired the calibration completely (log-loss
$2.0\to\sim 0.62$, into the tree regime) and lifted AUC $\sim$2 bands to
\num{0.7038} (Section~\ref{sec:anns-burn}) --- confirming the pathology was a
representation artifact, not an intrinsic property of nets. But the gap to the
trees only \emph{narrowed} (from $\sim$3.3 bands to $\sim$1), it did not close:
the net still lands below \texttt{catboost} and \texttt{xgboost}, the plain
\texttt{mlp} beat both the embedding and wide-deep topologies (entity embeddings
are starved at the corpus's low identity cardinality), and nothing cleared the
$\sim$0.71 save bar. The convergent finding stands; the standing lever for the
net, as for the trees, is identity \emph{resolution} --- raising the vocabulary
caps so the embedding topology can finally pay off.

\section{Source reports}
This synthesis consolidates six campaign reports from \texttt{experiments/}:
\begin{itemize}
  \item \texttt{2026-06-09-ae-xgboost-blend.md} --- autoencoder-latent +
    XGBoost blend on rotation (Section~\ref{sec:ae-blend}).
  \item \texttt{2026-06-09-svm-kernel-rotation.md} --- kernel SVM scan and
    conditional blend; first measured MAE noise band (Section~\ref{sec:svm}).
  \item \texttt{2026-06-09-rotation-classification.md} --- AUC bake-off, first
    measured AUC noise band, first registered rotation champion
    (Section~\ref{sec:auc}).
  \item \texttt{2026-06-12-classifier-bakeoff-p64.md} --- four-family
    classifier bake-off and xgboost depth/lr sweep on the new full-64-d corpus
    (Sections~\ref{sec:bakeoff-p64} and~\ref{sec:anns}).
  \item \texttt{2026-06-12-pyramid-mlp-arch-scan.md} --- six-geometry
    pyramid-MLP architecture scan, refuting the budget hypothesis
    (Sections~\ref{sec:mlp-scan} and~\ref{sec:anns}).
  \item \texttt{2026-06-12-burn-deep-topology.md} --- the Rust/\texttt{burn}
    \texttt{burn-deep-classifier} topology bake-off; resolves the
    representation-aware-net question, fixing calibration but not the AUC gap
    (Section~\ref{sec:anns-burn}).
\end{itemize}

\end{document}
